% Passes 5776--5783: q=5 Reye/Latin common-core projector and outer/sign split.
\subsection{The q=5 Reye and heavy shells share one rank-nine Latin core}
\label{sec:pass5776-5783-reye-latin-common-core}

The Pass~5667--5674 q=5 multidesign certificate contains two apparently
different incidence systems on the same moving twelve cover cells: the twelve
multiplicity-six heavy supports and the sixteen zero-containment triples forming
Reye's $12_4,16_3$ configuration.  Let
\[
 H\in\{0,1\}^{12\times12}
 \qquad\text{and}\qquad
 R\in\{0,1\}^{12\times16}
\]
be their point--block incidence matrices.  Exact arithmetic gives the stronger
identity
\[
 \boxed{
 HH^{T}-3J_{12}=RR^{T}-J_{12}=:K_9,
 }
\]
with
\[
 \boxed{K_9^2=4K_9},
 \qquad
 \boxed{\operatorname{rank}_{\mathbb Q}K_9=9}.
\]
Thus
\[
 \boxed{P_9=\frac14K_9}
\]
is a rational orthogonal projector canonically extracted from either shell.

The same matrix reconstructs an intrinsic imprimitivity system.  The matrix
\[
 A_{\rm blk}=3I_{12}-K_9
\]
is exactly $3K_4$.  In the frozen moving-twelve labels its components are
\[
 \{1,3,8,10\},\qquad
 \{2,5,7,12\},\qquad
 \{4,6,9,11\},
\]
and, after block-ordering,
\[
 K_9=I_3\otimes(4I_4-J_4),
 \qquad
 P_9=I_3\otimes\left(I_4-\frac14J_4\right).
\]
Hence the protected common nine-space is simply the direct sum of the three
within-block zero-sum spaces.

The sixteen Reye triples meet each of these four-point blocks once, and every
cross-pair from any two blocks occurs in exactly one triple.  Therefore the
zero shell is a transversal design
\[
 \boxed{TD(3,4)},
\]
equivalently a Latin square of order four.  With the deterministic component
ordering of the verifier its table is
\[
\begin{pmatrix}
2&0&3&1\\
0&2&1&3\\
3&1&2&0\\
1&3&0&2
\end{pmatrix},
\]
and the symbol relabel $(0,1,2,3)\mapsto(1,3,0,2)$ turns it into
\[
 \boxed{
 \begin{pmatrix}
 0&1&2&3\\
 1&0&3&2\\
 2&3&0&1\\
 3&2&1&0
 \end{pmatrix}
 }
\]
---the Cayley table of the Klein group $V_4\cong\mathbb F_2^2$ under XOR.
In the original thirteen-cover positions one explicit chart is
\[
\begin{aligned}
 r=0,1,2,3&\longleftrightarrow 2,8,11,3,\\
 c=0,1,2,3&\longleftrightarrow 13,6,5,9,\\
 s=0,1,2,3&\longleftrightarrow 7,10,12,4,
\end{aligned}
\qquad
\boxed{s=r\oplus c}.
\]
The row/column/symbol names remain autoparatopy gauge choices; the recovered
$TD(3,4)$ object is intrinsic.

The heavy shell is also recoverable from this Latin object alone.  Among all
\[
 \binom42^3=216
\]
balanced six-subsets choosing two points from each four-block, the number of
contained Reye triples has exact spectrum
\[
 \boxed{0^{12}\oplus2^{192}\oplus4^{12}}.
\]
The four-line class is exactly the twelve intercalate supports of the Klein
square, while the zero-line class is exactly the twelve multiplicity-six heavy
supports.  Complementation in the moving twelve exchanges the two exceptional
classes:
\[
 \boxed{
 \text{heavy support}
 =\text{complement of a Klein intercalate support}.
 }
\]
This reconstructs the heavy shell, but does not select a preferred bijection
between individual points and individual heavy supports.

The common projector resolves the earlier character overlap.  The three
transitive carrier actions have rank three and pairwise permutation-character
inner product two.  Since both incidence matrices have rank ten and their
centered parts have rank nine, their rational permutation modules decompose by
dimension as
\[
 \boxed{
 \mathbb Q^{12}_{P}=\mathbf1\oplus W_9\oplus U_{2,P},
 }
\]
\[
 \boxed{
 \mathbb Q^{12}_{H}=\mathbf1\oplus W_9\oplus U_{2,H},
 }
\]
and
\[
 \boxed{
 \mathbb Q^{16}_{L}=\mathbf1\oplus W_9\oplus V_6.
 }
\]
Thus every pair shares exactly the trivial line and the same irreducible
nine-space $W_9$.

This also makes the outer/sign distinction structural.  The order-two carrier
outer automorphism of Pass~5674 exchanges the point- and heavy-stabilizer
classes.  Because $W_9$ is their unique common nontrivial constituent, the
involution fixes $W_9$ and swaps the two-dimensional spokes:
\[
 \boxed{W_9\mapsto W_9,
 \qquad U_{2,P}\leftrightarrow U_{2,H}.}
\]
By contrast the point-side sign tensor twist is disjoint from the heavy and line
modules by Pass~5674.  Its frozen odd-element point fixed-count distribution
$0^{168},2^{108},6^{12}$ also gives
\[
 \langle\varepsilon\pi_P,\pi_P\rangle
 =\frac{3\cdot576-2(108\cdot2^2+12\cdot6^2)}{576}=0.
\]
Hence
\[
 \boxed{
 \varepsilon\otimes\mathbb Q^{12}_P
 \perp
 \mathbb Q^{12}_P\oplus\mathbb Q^{12}_H\oplus\mathbb Q^{16}_L.
 }
\]
The source outer involution preserves the common core and exchanges spokes;
the sign twist moves the whole point packet to a disjoint character sector.

Finally, this resolves a possible order-$576$ ambiguity.  Pass~5300 concerns
the projective-symplectic Hoffman-cover stabilizer
$2_+^{1+4}:(S_3\times C_3)$ with centre $C_2$.  Pass~5417 instead defines
\texttt{act} as the cover-cell permutation image of the full graph-automorphism
stabilizer; its frozen producer certificate is
\[
 \boxed{2^4:(S_3\times S_3),\qquad |\texttt{act}|=576,
 \qquad Z(\texttt{act})=1.}
\]
The present Reye/Latin theorem concerns this latter centreless group.  The two
order-$576$ groups are not identified.

\paragraph{Prior art and evidence boundary.}
The classical Reye configuration and its order-$576$ Levi automorphism group
are prior art.  In particular Monson--Pellicer--Williams, \emph{The tomotope},
Ars Math. Contemp. \textbf{5} (2012), 355--370,
DOI~10.26493/1855-3974.189.e64, identify the tomotope medial layer graph with
the Reye Levi graph.  The new statement here is the exact q=5 multidesign
bridge: identical centered Gram projector, intrinsic Klein-$V_4$ $TD(3,4)$,
heavy/intercalate complement class, and the resulting representation split.
No continuum, particle, gauge-force, mass/coupling, or physical-unification
claim is inferred.
